\documentclass[../main.tex]{subfiles}
\begin{document}

\section{Архитектура системы}

Система построена по клиент-серверной схеме с выделенным сервисом
сбора и хранения метаданных. Общая архитектура приведена
на рисунке~\ref{fig:architecture}.

\begin{figure}[H]
  \centering
  \resizebox{\linewidth}{!}{%
    \input{\subfix{diagrams/architecture}}%
  }
  \caption{Архитектура системы}
  \label{fig:architecture}
\end{figure}

В верхней части архитектуры расположен блок REPOSITORY, объединяющий
три независимых источника данных. Каждый источник работает как
отдельный контейнер с собственной изолированной сетью Docker, что
исключает возможность прямого сетевого взаимодействия между ними
и запрещает межбазовые SQL-запросы на уровне СУБД. Первый источник
(VM1) содержит базу \field{authorities\_dict} на PostgreSQL~16
и хранит медленно меняющиеся справочные данные: виды деятельности,
иерархический справочник полномочий, базовые должности, шаблоны
доверенностей и падежные формы. Второй источник (VM2) содержит базу
\field{authorities\_ops} на PostgreSQL~16 и хранит операционные
данные: организации, паспорта, сотрудников, должности
и доверенности. Третий источник (VM3) содержит архивную базу
\field{authorities\_archive} на MariaDB~11 с историческими версиями
доверенностей, журналом аудита и снимками справочников. Такое разделение источников по разным СУБД и разным
контейнерам позволяет проверить, что система одинаково работает
с разнородными источниками и не полагается на особенности
конкретной реализации.

Центральным элементом архитектуры является сервис SERVICE,
реализованный на Java~25 с использованием фреймворка Quarkus~3.
Сервис подключается к каждому источнику по JDBC в режиме <<только
чтение>> и собирает описания структуры данных: схемы, таблицы,
колонки, ограничения и индексы. Соединение с источником открывается
по требованию, читает метаданные одним пакетом запросов и сразу
закрывается. Это гарантирует, что сервис не вмешивается в работу
источника, не блокирует таблицы и не требует его остановки
или перестройки. Сбор метаданных для трёх источников может
выполняться параллельно, поскольку каждый источник работает
в собственной сети и не зависит от остальных.

Собранные описания структуры сохраняются в агрегирующей базе
данных AGGREGATOR\_DB, работающей на PostgreSQL~17. Она хранит
снапшоты схем всех источников в нормализованной форме: таблицы
с префиксом \field{sc\_} повторяют иерархию самих СУБД (инстанс,
схема, таблица, колонка, ограничение, индекс) и дополнительно
содержат журнал федеративных запросов. Ключи, ограничения
и индексы, в том числе составные,
хранятся реляционно: шапка объекта вынесена в одну таблицу, состав
колонок~--- в дочернюю, с сохранением порядка следования. Это
позволяет сравнивать описания двух снапшотов структурно,
а не по тексту, и корректно обрабатывать составные ключи любой
длины.

Взаимодействие между сервисом и агрегирующей базой данных
двунаправленное. Операция \field{INSERT} используется при
сохранении нового снапшота: сервис записывает собранные описания
в таблицы \field{sc\_*} в рамках одной транзакции, что исключает
появление частично сохранённых снапшотов. Операция \field{SELECT}
используется при чтении: когда сервису нужно отдать клиенту
сохранённый каталог, сравнить две схемы или построить план
федеративного запроса. Агрегирующая база данных не содержит
бизнес-данных источников~--- только описания их структуры. Такое
разделение позволяет хранилищу оставаться компактным и не зависеть
от объёма данных в самих источниках.

Клиентская часть (CLIENT) взаимодействует с сервисом по HTTP.
HTTP-запросы используются для получения списка источников, каталога
схем, результата сравнения двух каталогов и выполнения федеративных
SQL-запросов. Все операции возвращают полный JSON-ответ
и завершаются сразу после отправки данных. После загрузки схемы
клиент повторно запрашивает каталог, чтобы отобразить его актуальное
состояние.

\end{document}